% Options for packages loaded elsewhere
\PassOptionsToPackage{unicode}{hyperref}
\PassOptionsToPackage{hyphens}{url}
\PassOptionsToPackage{dvipsnames,svgnames,x11names}{xcolor}
\documentclass[
  spanish,
  11pt,
]{article}
\usepackage{xcolor}
\usepackage[margin=2.5cm]{geometry}
\usepackage{amsmath,amssymb}
\setcounter{secnumdepth}{-\maxdimen} % remove section numbering
\usepackage{iftex}
\ifPDFTeX
  \usepackage[T1]{fontenc}
  \usepackage[utf8]{inputenc}
  \usepackage{textcomp} % provide euro and other symbols
\else % if luatex or xetex
  \usepackage{unicode-math} % this also loads fontspec
  \defaultfontfeatures{Scale=MatchLowercase}
  \defaultfontfeatures[\rmfamily]{Ligatures=TeX,Scale=1}
\fi
\usepackage{lmodern}
\ifPDFTeX\else
  % xetex/luatex font selection
\fi
% Use upquote if available, for straight quotes in verbatim environments
\IfFileExists{upquote.sty}{\usepackage{upquote}}{}
\IfFileExists{microtype.sty}{% use microtype if available
  \usepackage[]{microtype}
  \UseMicrotypeSet[protrusion]{basicmath} % disable protrusion for tt fonts
}{}
\makeatletter
\@ifundefined{KOMAClassName}{% if non-KOMA class
  \IfFileExists{parskip.sty}{%
    \usepackage{parskip}
  }{% else
    \setlength{\parindent}{0pt}
    \setlength{\parskip}{6pt plus 2pt minus 1pt}}
}{% if KOMA class
  \KOMAoptions{parskip=half}}
\makeatother
\ifLuaTeX
\usepackage[bidi=basic,shorthands=off]{babel}
\else
\usepackage[bidi=default,shorthands=off]{babel}
\fi
\ifLuaTeX
  \usepackage{selnolig} % disable illegal ligatures
\fi
\setlength{\emergencystretch}{3em} % prevent overfull lines
\providecommand{\tightlist}{%
  \setlength{\itemsep}{0pt}\setlength{\parskip}{0pt}}
% Símbolos Unicode (griegos, sub/superíndices, operadores) que Latin Modern no trae
\usepackage{newunicodechar}
\newfontfamily\simbolos{DejaVu Serif}[Scale=MatchLowercase]
\newunicodechar{α}{{\simbolos α}}
\newunicodechar{β}{{\simbolos β}}
\newunicodechar{δ}{{\simbolos δ}}
\newunicodechar{η}{{\simbolos η}}
\newunicodechar{π}{{\simbolos π}}
\newunicodechar{ρ}{{\simbolos ρ}}
\newunicodechar{τ}{{\simbolos τ}}
\newunicodechar{χ}{{\simbolos χ}}
\newunicodechar{ᵢ}{{\simbolos ᵢ}}
\newunicodechar{⁰}{{\simbolos ⁰}}
\newunicodechar{⁶}{{\simbolos ⁶}}
\newunicodechar{⁷}{{\simbolos ⁷}}
\newunicodechar{⁹}{{\simbolos ⁹}}
\newunicodechar{⁻}{{\simbolos ⁻}}
\newunicodechar{ⁿ}{{\simbolos ⁿ}}
\newunicodechar{₀}{{\simbolos ₀}}
\newunicodechar{₁}{{\simbolos ₁}}
\newunicodechar{₂}{{\simbolos ₂}}
\newunicodechar{₊}{{\simbolos ₊}}
\newunicodechar{ₖ}{{\simbolos ₖ}}
\newunicodechar{ₙ}{{\simbolos ₙ}}
\newunicodechar{∈}{{\simbolos ∈}}
\newunicodechar{∉}{{\simbolos ∉}}
\newunicodechar{∝}{{\simbolos ∝}}
\newunicodechar{≈}{{\simbolos ≈}}
\newunicodechar{≠}{{\simbolos ≠}}
\newunicodechar{≤}{{\simbolos ≤}}
\newunicodechar{≥}{{\simbolos ≥}}
\newunicodechar{⊆}{{\simbolos ⊆}}
\usepackage{bookmark}
\IfFileExists{xurl.sty}{\usepackage{xurl}}{} % add URL line breaks if available
\urlstyle{same}
\hypersetup{
  pdftitle={Instancias},
  pdflang={es},
  colorlinks=true,
  linkcolor={Maroon},
  filecolor={Maroon},
  citecolor={Blue},
  urlcolor={Blue},
  pdfcreator={LaTeX via pandoc}}

\title{Instancias}
\author{}
\date{}

\begin{document}
\maketitle

\section{Origen y verificación}\label{origen-y-verificaciuxf3n}

\begin{itemize}
\tightlist
\item
  \texttt{envelope6.hcp}: ejemplo oficial de 6 vértices de la página del
  FHCP (3 ciclos hamiltonianos conocidos).
\item
  \texttt{graph1,\ graph2,\ graph3} (depuración, n=66/70/78) y
  \texttt{graph24,\ graph41,\ graph172} (n=198/300/1002): tomados de una
  \textbf{copia pública} del FHCP Challenge Set en GitHub (repositorio
  cgungor/ham-path, carpeta FHCPCS), no del sitio oficial de Flinders.
\item
  Comprobaciones hechas: la copia tiene 1001 grafos, n entre 66 y 9528 y
  media de \textasciitilde3098 vértices, que coincide con lo que publica
  Flinders; los 1001 cargan sin errores; encabezado ``FHCP Challenge
  Set. http://fhcp.edu.au/fhcpcs''.
\item
  \textbf{Pendiente:} compararlos con los archivos oficiales (FHCPCS.7z)
  si logras descargarlos (comparar \texttt{sha256} del catálogo contra
  el de los archivos oficiales). Mientras tanto, declarar en el informe
  que la fuente fue una copia pública.
\item
  \texttt{ejemplo6.hcp} es un grafo propio de prueba, no es del FHCP.
\end{itemize}

\section{Catálogo}\label{catuxe1logo}

\texttt{data/catalogo\_fhcp.csv} lista las 1001 instancias (archivo, n,
m, grado mín/máx, sha256). Se regenera con:
\texttt{python\ -m\ experiments.catalogo\ CARPETA\_CON\_HCP\ salida.csv}

\end{document}
